% ============================================================================
% Chapter 14 --- Orchestration plane
% ----------------------------------------------------------------------------
% Purpose : the cross-layer orchestrator that issues commands following a
%           TAC activation. The architecture's principal acknowledged gap
%           (D5.3): no single standard. Belongs to the Control and
%           Governance plane (D5.1).
% Page target: ~5 pages.
% Dependencies: Ch. 5 (D5.3 declares this gap), every layer that the
%               orchestrator touches (Ch. 6--12), Ch. 12 (events through
%               canonical schema), Ch. 16 (conformance for orchestration).
% Mirrors layer chapter template established in Ch. 6 (D6.5).
% ----------------------------------------------------------------------------
% Decisions registered in _coherence-EN.md:
%   D14.1 four required properties of the orchestration plane (the contract)
%   D14.2 reference instantiation = Argo Workflows + Crossplane + GitOps
%   D14.3 alternatives = ServiceNow/Ansible Tower; custom on Temporal
%   D14.4 the gap is bounded by composition discipline and declarative state
%   D14.5 orchestration sovereignty scores
% ============================================================================

\chapter{Orchestration plane}
\label{ch:orchestration}

\begin{caveatbox}[Dependencies and acknowledged gap]
\noindent This chapter applies the layer template of
\trchap{ch:identity}{} to the orchestration plane. As
\trchap{ch:architectural-principles}{} (\S\ref{sec:ap-standards}, D5.3) acknowledges,
this is the layer where no single standard consolidates the
required functionality. The chapter's purpose is to bound the
consequences of that gap by composing components on standards-aligned
sub-interfaces, by keeping the composition declarative, and by
treating the orchestrator itself as replaceable infrastructure
rather than as institutional knowledge in code. The layer belongs
to the \emph{Control and governance plane}.
\end{caveatbox}

The orchestration plane is the glue that turns a TAC activation
decision into the cross-layer commands that realise it. When the
policy layer returns ``allow role X for resource Y in zone Z'' for
a given subject, the orchestration plane is what tells the network
to attach the subject to zone Z, tells the compute layer to
provision the workload context, tells the storage layer to bind the
appropriate volumes, and tells the observability layer to start the
correlation trace. The plane is consequently both critical (without
it nothing happens) and exposed to lock-in (the supplier whose
orchestration patterns the institution adopts becomes a structural
dependency).

%-----------------------------------------------------------------------------
\section{Functional role and interface contract}
\label{sec:or-contract}

\begin{contractbox}[Orchestration plane]
\noindent A compliant orchestration plane, in the sense of this
architecture, must exhibit four properties rather than speak a
single standard:

\begin{itemize}[leftmargin=1.5em,nosep]
  \item \textbf{Event-driven.} The plane reacts to TAC activation
    and deactivation events emitted by the policy layer; it is not
    polled, scheduled, or driven by side channels.
  \item \textbf{Idempotent.} Re-issuing the same activation
    produces the same end state; partial failures and retries are
    safe by construction, not by good operational fortune.
  \item \textbf{Observable.} Every step the plane takes is recorded
    as an event conformant to the canonical TAC schema of
    \trchap{ch:observability}{} (\S\ref{sec:ob-schema}), with the
    \texttt{correlation\_id} preserved across the entire activation.
  \item \textbf{Reversible.} The deactivation of a TAC undoes
    every consequential change the plane made under it, in a
    documented procedure that the institution can audit and exercise.
\end{itemize}
\end{contractbox}

\paragraph{Why properties rather than a standard.} The four
properties together form the plane's semantic contract. They can
be satisfied by different technical compositions: workflow engines
with explicit step-by-step execution, declarative reconcilers that
converge on a target state, or hybrid arrangements that combine
both. The architecture commits to the properties; the choice of
composition belongs to the institution.

%-----------------------------------------------------------------------------
\section{Reference instantiation: declarative GitOps reconciliation with composable workflows}
\label{sec:or-reference}

\begin{instantiationbox}[Orchestration plane]
\noindent\textbf{Topology.} A GitOps repository expresses the
declarative target state of the institution's composed
infrastructure --- network attachment policies, compute namespaces,
storage bindings, observability collectors, identity client
registrations. A declarative reconciler --- of which open-source
continuous-delivery controllers such as Argo~CD~\cite{argocd} or
Flux are representative --- ensures that the running state converges
on the repository state. Workflow logic that does not naturally fit
the declarative model --- for example, the multi-step provisioning
of a research workspace that needs explicit ordering and conditional
branches --- is expressed in a container-native workflow engine, of
which Argo~Workflows~\cite{argo-workflows} is a representative
open-source instance, with workflow definitions also held in the
GitOps repository. A control-plane provider framework --- the
open-source Crossplane~\cite{crossplane} project being one such
instance --- provides the infrastructure-as-code layer that lets the
same declarative model reach beyond Kubernetes itself, into the
virtualisation platform, into the network layer's RADIUS
configuration, and into the storage layer's volume management. TAC
activation events from the policy layer are received as
CloudEvents and trigger either a reconciliation or a workflow,
depending on the activation type. Every step emits events
conformant to the canonical schema.
\end{instantiationbox}

\paragraph{Operational considerations.} The reference instantiation
is the most demanding of the layer chapters in terms of
\emph{integration} skills, although each individual component is
broadly skilled. The principal recurring effort is composition:
keeping the GitOps repository, the workflow corpus, and the
control-plane provider definitions consistent with each other and
with the upstream layers they touch. An indicative envelope, for \loc{institution-type}, would be on the order of two to four engineers shared
with the platform-engineering function. The labour market is broad
for the individual components and currently narrow for their
integrated composition.

\paragraph{Composition discipline.} The architectural posture is to
treat the composition itself as institutional artefact: the GitOps
repository, the workflow corpus, and the provider configurations
are versioned and reviewed with the same discipline as the policy
corpus of \trchap{ch:policy}{}. The institution that lets these
artefacts drift loses the property of replaceability for the
orchestration plane, even if each individual component remains
replaceable.

\begin{realworldbox}[Orchestration patterns in industrial and research settings]
\noindent The cloud-native orchestration ecosystem --- the family of
open-source GitOps reconcilers, container-native workflow engines,
control-plane provider frameworks, pipeline engines, and durable
workflow engines, of which Argo, Flux, Crossplane, Tekton, and
Temporal are representative members --- is broadly deployed in
industrial settings, with publicly available reference architectures
from the projects' member organisations.
Specific academic deployments tend to combine these patterns with
institution-specific composition, and are less extensively
documented in public fora. The phase-2 international university
network anticipated by this Reference would be one of the contexts
in which a shared library of academic orchestration compositions
might emerge.
\end{realworldbox}

%-----------------------------------------------------------------------------
\section{Alternative~1 --- ITSM-driven orchestration with Ansible}
\label{sec:or-alt1}

\begin{alternativebox}[ServiceNow/Jira + Ansible]
\noindent The institution operates an existing IT service
management platform (ServiceNow, Atlassian Jira Service Management,
or comparable) that triggers Ansible playbooks (in Ansible Automation
Platform or open Ansible) when service requests match TAC
activation patterns. The four contractual properties are realised
through workflows defined in the ITSM platform and idempotent
playbooks executed by Ansible.
\end{alternativebox}

\paragraph{Where it adds value.} For institutions with a mature
ITSM platform and existing investment in Ansible operations, the
operational continuity is substantial. The audit trail produced by
the ITSM platform is well-developed, and the integration with
change-management processes is direct.

\paragraph{Where it constrains the institution.} ITSM workflow
languages tend to be supplier-specific; expressing the four
properties in those languages typically produces logic that is not
portable across ITSM platforms. The architectural posture is to
keep the workflow logic minimal in the ITSM platform and to delegate
the consequential actions to standard Ansible playbooks held under
version control --- treating the ITSM platform as a
trigger-and-audit surface rather than as the orchestration logic
itself.

\paragraph{When it could be the right choice.} An institution
whose existing ITSM and Ansible operations carry the operational
weight, and whose sovereignty posture accepts the trade-offs in
exchange for the integration with established change-management
practice.

%-----------------------------------------------------------------------------
\section{Alternative~2 --- Custom orchestrator on a durable workflow engine}
\label{sec:or-alt2}

\begin{alternativebox}[Custom on Temporal or comparable]
\noindent The institution writes a bespoke orchestrator on a
durable workflow engine (Temporal.io~\cite{temporal} being the
principal open-source example, with comparable patterns available
in other engines). Workflows are expressed in a programming language with
the engine's SDK, with the engine providing durability,
replayability, and at-least-once execution.
\end{alternativebox}

\paragraph{Where it adds value.} For institutions whose
orchestration needs include long-running workflows (multi-day
provisioning, multi-stage research-environment lifecycle) where the
declarative model would be awkward, a durable workflow engine
provides the right semantics. Replay-based debugging and
deterministic re-execution are operationally substantial benefits.

\paragraph{Where it constrains the institution.} A custom
orchestrator written in a programming language risks accumulating
institutional knowledge in code that the standard interfaces do not
capture: an exit from the engine then becomes a re-implementation
of the workflows. The mitigation is to keep the workflows
expressing the four contractual properties and nothing else,
delegating consequential actions to the layer-specific standard
APIs.

\paragraph{When it could be the right choice.} An institution
with strong software-engineering capacity, with workflow needs that
genuinely exceed what declarative reconciliation expresses, and
with the discipline to keep the workflow corpus deliberately thin.

%-----------------------------------------------------------------------------
\section{Sovereignty grid applied}
\label{sec:or-sovereignty}

\begin{table}[H]
\centering\small
\caption{Per-dimension sovereignty profile of the orchestration
instantiations.}
\label{tab:or-sovereignty}
\begin{tabularx}{\textwidth}{%
  >{\hsize=0.5\hsize\linewidth=\hsize\bfseries\raggedright\arraybackslash}X%
  >{\hsize=1.2\hsize\linewidth=\hsize\centering\arraybackslash}X%
  >{\hsize=1.1\hsize\linewidth=\hsize\centering\arraybackslash}X%
  >{\hsize=1.2\hsize\linewidth=\hsize\centering\arraybackslash}X%
}
\toprule
Dimension & Reference (GitOps + control-plane providers) & ITSM + config-management & Custom durable-workflow engine \\
\midrule
Data          & 3 Robust      & 2 Established & 3 Robust       \\
Technical     & 3 Robust      & 1--2 Partial  & 2 Established  \\
Financial     & 2 Established & 1 Partial     & 2 Established  \\
Operational   & 2 Established & 3 Robust      & 1--2 Partial   \\
Institutional & 3 Robust      & 1 Partial     & 2 Established  \\
\bottomrule
\end{tabularx}
\end{table}

\noindent The orchestration plane is the layer where the
\emph{composition discipline} is more consequential than any
individual component choice. An institution that adopts the
reference composition without disciplined version control and
review can score worse, on the technical and institutional
dimensions, than an institution that adopts the ITSM alternative
with rigorous practice. The grid scores above assume disciplined
operation in each instantiation.

%-----------------------------------------------------------------------------
\section{Regulatory anchoring}
\label{sec:or-regulatory}

\noindent The orchestration plane is the deployment-and-change
substrate. It contributes to \gls{gdpr}~Art.~25 (data protection
by design through declarative defaults) and to NIS2~Art.~21(2)(c)
(business continuity through declarative redeployability), (e)
(security in development and maintenance), and (f) (effectiveness
assessment, through GitOps-driven reconciliation that exposes
drift). The layer operationalises 27001~A.5.30 (ICT readiness for
business continuity), A.8.32 (change management) and contributes
to NIST CSF \textsf{PROTECT} (PR.PS) and \textsf{RECOVER} (RC.RP
--- recovery plan execution). It also operationalises
COBIT~BAI06 (managed IT changes) and BAI10 (managed
configuration). Detailed mapping in
\trchap{ch:regulatory-mapping}{} \S\ref{sec:rm-gdpr},
\S\ref{sec:rm-nis2}, \S\ref{sec:rm-iso},
\S\ref{sec:rm-nist-csf}, \S\ref{sec:rm-cobit}.

%-----------------------------------------------------------------------------
\section{Common pitfalls}
\label{sec:or-pitfalls}

\begin{pitfallbox}[Orchestrator becomes the single point of failure]
\noindent An orchestration plane that holds critical state in its
runtime memory, in undocumented configuration, or in product-specific
formats becomes the SPOF for the entire architecture: when it
fails, no TAC can activate, deactivate, or be modified. The
mitigation is the GitOps discipline: every consequential decision
expressed in the repository, the runtime treated as derivable from
the repository, the recovery procedure exercised periodically.
\end{pitfallbox}

\begin{pitfallbox}[State not versioned]
\noindent Orchestration logic exists in three layers: the
declarative target state, the workflow definitions, and the
runtime configuration of the engines themselves. When the third
layer is not under version control, the orchestration plane has
no rebuildable definition: only its current running state. The
mitigation is to bring all three layers under version control,
including the engines' own configuration, and to treat
non-repository changes as incidents.
\end{pitfallbox}

\begin{pitfallbox}[Hidden coupling to a supplier SDK]
\noindent A workflow corpus written against a supplier's SDK
accumulates idioms that the SDK encourages but that the underlying
contract does not require. Migration to another engine becomes a
re-write rather than a re-host. The mitigation is to express
workflows against the four contractual properties of
\S\ref{sec:or-contract}, using the SDK only for the engine
interaction, and to maintain a small set of representative
workflows that have been ported between engines as part of the
exit-rehearsal of \trchap{ch:conformance}{}.
\end{pitfallbox}

\begin{pitfallbox}[Side-effects untracked]
\noindent A workflow that, in the course of provisioning a
resource, also creates side-channel changes (firewall rules added
manually, DNS records inserted by an out-of-band procedure,
secrets generated and stored in a non-canonical store)
accumulates technical debt that the declarative model does not
capture. The deactivation procedure then leaves orphans behind.
The mitigation is to require every change to flow through the
declarative model or through the workflow corpus --- there is no
admissible third path.
\end{pitfallbox}
